% ============================================================
\chapter{Modul 15: Brücke zu ROS2}
% ============================================================

Alles, was du in diesem Buch gelernt hast, führt hierher. ROS2 (Robot Operating System 2) ist das Standard-Framework für professionelle Robotik weltweit. Es strukturiert Roboterprogramme als Netzwerk kommunizierender \textbf{Nodes}.

\section{Lektion 41: Was ist ROS2?}

\begin{center}
\begin{tabular}{ll}
\toprule
\textbf{ROS2-Konzept} & \textbf{Python-Äquivalent (aus diesem Buch)} \\
\midrule
Node         & Klasse (Modul 9) \\
Topic        & Kommunikationskanal zwischen Nodes \\
Message      & Dictionary mit definiertem Schema (Modul 6) \\
Publisher    & Funktion, die Daten sendet \\
Subscriber   & Callback-Funktion, die Daten empfängt \\
Timer        & \texttt{while}-Schleife mit Taktung (Modul 3) \\
Parameter    & Konfigurationsdatei / Dictionary (Modul 7) \\
\bottomrule
\end{tabular}
\end{center}

\section{Lektionen 42--44: Python in ROS2}

\subsection*{Installation (Ubuntu 24.04)}

\begin{lstlisting}
# Terminal:
# sudo apt install ros-jazzy-desktop
# source /opt/ros/jazzy/setup.bash
\end{lstlisting}

\subsection*{Publisher -- Sensorwerte senden}

\begin{lstlisting}
import rclpy
from rclpy.node import Node
from std_msgs.msg import Float32
from sensor_msgs.msg import Range

class UltraschallPublisher(Node):
    """Liest Ultraschall-Sensor und veroeffentlicht den Wert."""

    def __init__(self):
        super().__init__("ultraschall_node")

        # Parameter (entspricht Konfigurationsdatei aus Modul 7)
        self.declare_parameter("min_abstand", 0.02)
        self.declare_parameter("max_abstand", 4.00)
        self.declare_parameter("frequenz_hz", 10.0)

        self.min_a = self.get_parameter("min_abstand").value
        self.max_a = self.get_parameter("max_abstand").value
        hz = self.get_parameter("frequenz_hz").value

        # Publisher (sendet auf Topic /abstand)
        self.publisher = self.create_publisher(Float32, "/abstand", 10)

        # Timer ersetzt die while-Schleife aus Modul 3
        dt = 1.0 / hz
        self.timer = self.create_timer(dt, self.timer_callback)

        self.get_logger().info(f"Ultraschall-Node gestartet ({hz} Hz)")

    def timer_callback(self):
        """Wird automatisch mit der konfigurierten Frequenz aufgerufen."""
        # Hier: echten Sensor lesen statt Zufallswert
        import random
        wert = round(random.uniform(self.min_a, self.max_a), 3)

        msg = Float32()
        msg.data = wert
        self.publisher.publish(msg)
        self.get_logger().debug(f"Abstand: {wert} m")


def main():
    rclpy.init()
    node = UltraschallPublisher()
    try:
        rclpy.spin(node)       # = while True: bearbeite_events()
    except KeyboardInterrupt:
        pass
    finally:
        node.destroy_node()
        rclpy.shutdown()

if __name__ == "__main__":
    main()
\end{lstlisting}

\subsection*{Subscriber -- Sensorwerte empfangen}

\begin{lstlisting}
import rclpy
from rclpy.node import Node
from std_msgs.msg import Float32
from geometry_msgs.msg import Twist

class HindernisvermeidungNode(Node):
    """Empfaengt Abstandswerte und steuert Bewegung."""

    def __init__(self):
        super().__init__("hindernis_node")
        self.declare_parameter("stopp_abstand", 0.25)
        self.stopp_a = self.get_parameter("stopp_abstand").value

        # Subscriber: lauscht auf /abstand
        self.sub = self.create_subscription(
            Float32, "/abstand", self.abstand_callback, 10)

        # Publisher: sendet Bewegungsbefehle
        self.cmd_pub = self.create_publisher(Twist, "/cmd_vel", 10)

        self.get_logger().info("Hindernisvermeidung aktiv.")

    def abstand_callback(self, msg):
        """Wird aufgerufen, wenn eine neue Abstandsmessung eintrifft."""
        abstand = msg.data
        twist = Twist()

        # Entscheidungslogik (aus Modul 2)
        if abstand < self.stopp_a:
            twist.linear.x  = 0.0
            twist.angular.z = 0.5   # Drehen
            self.get_logger().warn(f"Hindernis bei {abstand:.2f} m!")
        else:
            twist.linear.x  = 0.3
            twist.angular.z = 0.0

        self.cmd_pub.publish(twist)
\end{lstlisting}

\section{Lektion 45: Dein erster mobiler Roboter}

\begin{lstlisting}
# Vollstaendiges minimales ROS2-Paket:
# Struktur:
#   mein_roboter/
#   ├── setup.py
#   ├── package.xml
#   └── mein_roboter/
#       ├── __init__.py
#       ├── ultraschall_node.py   (Publisher aus Lektion 42)
#       └── hindernis_node.py     (Subscriber aus Lektion 43)

# Terminal-Befehle:
# colcon build --packages-select mein_roboter
# source install/setup.bash
# ros2 run mein_roboter ultraschall_node &
# ros2 run mein_roboter hindernis_node
\end{lstlisting}

\begin{merksatz}
Die komplette Python-Grundlage aus diesem Buch steckt in jedem ROS2-Node: Klassen (Modul 9), Fehlerbehandlung (Modul 10), Funktionen (Modul 4), Dictionaries für Parameter (Modul 6), Schleifen als Timer (Modul 3). ROS2 ist kein neues Konzept -- es ist eine strukturierte Anwendung aller bisherigen Konzepte.
\end{merksatz}

\begin{praxis}
Mit diesem Wissen kannst du:
\begin{itemize}
    \item Sensordaten aus ROS2-Topics lesen und verarbeiten
    \item Bewegungsbefehle an reale Roboter senden (\texttt{/cmd\_vel})
    \item Eigene ROS2-Nodes für deinen MobileRobot-1 schreiben
    \item Den Roboter mit deinem ML-Modell aus Modul 14 steuern
\end{itemize}
\end{praxis}

\begin{uebung}[title=Abschlussprojekt: Autonomer Patrouillier-Roboter]
Entwerfe (auf Papier oder als Code-Skeleton) einen vollständigen ROS2-Stack für einen autonomen Patrouillier-Roboter:
\begin{enumerate}
    \item \textbf{sensor\_node}: Liest Ultraschall und IMU, publiziert auf \texttt{/scan} und \texttt{/imu}
    \item \textbf{navigation\_node}: Subscribt auf \texttt{/scan}, berechnet Weg, publiziert auf \texttt{/cmd\_vel}
    \item \textbf{monitor\_node}: Subscribt auf alle Topics, loggt Daten, erkennt Anomalien
    \item \textbf{Parameter}: Alle Schwellenwerte als ROS2-Parameter (kein Hardcoding)
\end{enumerate}
Schreibe für jeden Node den \texttt{\_\_init\_\_}-Konstruktor und die wichtigste Callback-Funktion.
\end{uebung}
